% ============================================================================
% TCC — Predição de Direção de Preços de Ações com NLP e Sentimento de Notícias
% Formato: abnTeX2 (ABNT NBR 14724:2011)
% Compilar: pdflatex tcc.tex && bibtex tcc && pdflatex tcc.tex && pdflatex tcc.tex
% ============================================================================
\documentclass[
    12pt,
    openright,
    oneside,
    a4paper,
    english,
    brazil
]{abntex2}

% --- Pacotes ---
\usepackage[utf8]{inputenc}
\usepackage[T1]{fontenc}
\usepackage{amsmath,amssymb}
\usepackage{mathptmx} % Times New Roman (texto e matemática) --- exigência ABNT/PUC-SP
\usepackage{graphicx}
\usepackage[a4paper,top=3cm,left=3cm,bottom=2cm,right=2cm]{geometry} % margens ABNT NBR 14724
\usepackage{booktabs}
\usepackage{tabularx}
\usepackage{multirow}
\usepackage{float}
\usepackage{longtable}
\usepackage[expansion=false]{microtype} % Times (psnfss) não permite expansão automática
\usepackage{hyperref}
\usepackage[brazilian,hyperpageref]{backref}
\usepackage[alf]{abntex2cite}
\hypersetup{hidelinks} % links sem caixa colorida (versão banca/impressa)
\usepackage{xcolor}
\usepackage{listings}
\usepackage{enumitem}

% --- Fonte única (Times) em todo o trabalho, inclusive títulos ---
% PUC-SP exige fonte única (Times New Roman ou Arial); abnTeX2 usa sans nos
% títulos por padrão --- aqui forçamos serif (Times) também nos títulos.
\renewcommand{\ABNTEXchapterfont}{\rmfamily}
\renewcommand{\ABNTEXsectionfont}{\rmfamily}
\renewcommand{\ABNTEXsubsectionfont}{\rmfamily}
\renewcommand{\ABNTEXsubsubsectionfont}{\rmfamily}

% --- Caminho das figuras ---
\graphicspath{{../9.baselines/}{figuras/}}

% --- Configurações de listings ---
\lstset{
    basicstyle=\ttfamily\small,
    breaklines=true,
    frame=single,
    numbers=left,
    numberstyle=\tiny\color{gray},
    keywordstyle=\color{blue},
    commentstyle=\color{green!50!black},
    stringstyle=\color{red!60!black},
    language=Python,
}

% --- Configurações de backref ---
\renewcommand{\backrefpagesname}{Citado na(s) página(s):~}
\renewcommand{\backref}{}
\renewcommand*{\backrefalt}[4]{
    \ifcase #1
        Nenhuma citação no texto.
    \or
        Citado na página #2.
    \else
        Citado #1 vezes nas páginas #2.
    \fi
}

% --- Informações do documento ---
\titulo{Predição de Direção de Preços de Ações Brasileiras\\com Sentimento de Notícias Financeiras:\\Pipeline, Ilusão e Autocorreção Metodológica}
\autor{André Takeo Loschner Fujiwara}
\local{São Paulo}
\data{2026}
\orientador{Prof.\ Eric Bacconi Gonçalves}
\instituicao{%
    PUC-SP
    \par Curso de Ciência de Dados e Inteligência Artificial
}
\tipotrabalho{Trabalho de Conclusão de Curso}
\preambulo{Trabalho de Conclusão de Curso apresentado ao Curso de Ciência de Dados e Inteligência Artificial da PUC-SP como requisito parcial para a obtenção do grau de Bacharel.}

% --- Compila o índice ---
\makeindex

